% TOP_PROBLEM: {"id": 30004859, "title": "Zariski's Multiplicity Conjecture", "category_id": 5, "category": {"id": 5, "name": "algebraic_geometry", "display_name": "Algebraic Geometry", "description": "Geometric objects defined by polynomial equations.", "slug": "algebraic-geometry", "order_index": 5, "created_at": "2026-07-31T15:26:25.671Z"}, "status": "open"}
% FIRST_POSTED: 2026-09-25
% ENTRY_KIND: statement_only
% !TeX program = lualatex
% Definition and sources only; this file is not an LLM solution attempt.
\documentclass[11pt]{article}
\usepackage[margin=25mm]{geometry}
\usepackage{fontspec}
\setmainfont{Latin Modern Roman}
\usepackage{amsmath,amssymb,mathtools}
\usepackage{unicode-math}
\setmathfont{Latin Modern Math}
\usepackage{xurl,hyperref}
\hypersetup{colorlinks=true,urlcolor=blue}
\setcounter{secnumdepth}{0}
\setlength{\parindent}{0pt}
\setlength{\parskip}{5pt}
\setlength{\emergencystretch}{3em}
\begin{document}
\section*{\texorpdfstring{Zariski's Multiplicity Conjecture}{Zariski's Multiplicity Conjecture}}\label{30004859}
\subsection{Definitions and mathematical statement}
Let $f,g:({\mathbb{C}}^n,0)\to({\mathbb{C}},0)$ be nonzero reduced holomorphic germs. Their multiplicities are their orders of vanishing, namely the least total degrees of nonzero terms in their Taylor expansions. Embedded topological equivalence means there is a germ of a homeomorphism of ambient neighborhoods, fixing zero, that maps $V(f)$ onto $V(g)$. The conjecture is
\[
 ({\mathbb{C}}^n,V(f),0)\cong_{\mathrm{top}}({\mathbb{C}}^n,V(g),0)
 \quad\Longrightarrow\quad\operatorname{mult}_0f=\operatorname{mult}_0g.
\]
Reducedness is essential because taking a power preserves the zero set while changing the order of vanishing. No isolated-singularity or quasihomogeneity hypothesis is included in the full target.
\par\smallskip\textit{Formulation and concept references:} \href{https://arxiv.org/html/2309.02624}{[S1]}.

\subsection{Short English statement}
Does the embedded topological shape of a reduced complex hypersurface singularity determine its multiplicity?
\subsection{Sources}
\begin{itemize}
\item[S1]Zariski's multiplicity conjecture for quasihomogeneous hypersurfaces with non-isolated singularities. Accessed 2026-09-01.
\url{https://arxiv.org/html/2309.02624}.
\textit{Formulation source.}
\end{itemize}
\end{document}
